\documentclass[a4paper,12pt]{article}

% ------------------------------------------------------------
% Кодировка и язык
% ------------------------------------------------------------
\usepackage[T2A]{fontenc}
\usepackage[utf8]{inputenc}
\usepackage[russian]{babel}

% ------------------------------------------------------------
% Математика
% ------------------------------------------------------------
\usepackage{amsmath,amssymb,amsthm,mathtools}
\usepackage{icomma}

% ------------------------------------------------------------
% Шрифт и типографика
% ------------------------------------------------------------
\usepackage{helvet}
\renewcommand{\familydefault}{\sfdefault}
\usepackage{microtype}
\usepackage{setspace}
\onehalfspacing
\usepackage{parskip}

% ------------------------------------------------------------
% Страница
% ------------------------------------------------------------
\usepackage[
  left=20mm,
  right=20mm,
  top=20mm,
  bottom=22mm
]{geometry}

% ------------------------------------------------------------
% Графика
% ------------------------------------------------------------
\usepackage{tikz}
\usetikzlibrary{calc}
\usepackage{tkz-euclide}
\usepackage{pgfplots}
\pgfplotsset{compat=1.18}

% ------------------------------------------------------------
% Таблицы и списки
% ------------------------------------------------------------
\usepackage{tabularx,booktabs,longtable}
\usepackage{enumitem}

% ------------------------------------------------------------
% Служебные пакеты
% ------------------------------------------------------------
\usepackage{xcolor}
\usepackage{fancyhdr}
\usepackage{hyperref}
\hypersetup{hidelinks}

% ------------------------------------------------------------
% Цветовая палитра
% ------------------------------------------------------------
\definecolor{accent}{HTML}{008080}
\definecolor{darkaccent}{HTML}{004D4D}
\definecolor{textgray}{HTML}{333333}
\definecolor{softgray}{HTML}{6B7280}
\definecolor{lightaccent}{HTML}{E8F3F3}

\color{textgray}

% ------------------------------------------------------------
% Заголовки и колонтитулы
% ------------------------------------------------------------
\pagestyle{fancy}
\fancyhf{}
\fancyhead[L]{\small\color{softgray}График координаты и линейная функция}
\fancyhead[R]{\small\color{softgray}Володя Хачатурян}
\fancyfoot[C]{\small\thepage}
\renewcommand{\headrulewidth}{0.3pt}
\renewcommand{\headrule}{\hbox to\headwidth{%
  \color{accent}\leaders\hrule height \headrulewidth\hfill}}

\setcounter{secnumdepth}{0}

\newcommand{\sectionline}{%
  \par\noindent\color{accent}\rule{\linewidth}{0.5pt}\color{textgray}\par
}

\newcommand{\important}[1]{\textbf{\color{darkaccent}#1}}

\newcommand{\checkitem}{\(\square\)\ }

\newcommand{\daytitle}[1]{%
  \vspace{0.7em}
  {\large\bfseries\color{darkaccent}#1}\par
  \vspace{0.2em}
  \noindent\color{accent}\rule{\linewidth}{0.4pt}\color{textgray}
  \vspace{0.3em}
}

\begin{document}

% ============================================================
% ТИТУЛЬНЫЙ ЛИСТ
% ============================================================

\begin{titlepage}
\thispagestyle{empty}

\begin{center}

\vspace*{2.2cm}

{\Large\color{softgray}\textbf{ЧЕК-ЛИСТ}}

\vspace{0.7cm}

{\Huge\bfseries\color{darkaccent}
График координаты\\[0.25cm]
при равномерном движении
}

\vspace{0.45cm}

{\LARGE\bfseries
Линейная функция и определение\\
параметров движения по графику
}

\vspace{1.2cm}

{\large
Подготовка к ОГЭ по физике\\
с необходимым повторением алгебры
}

\vfill

{\large
\textbf{Лёвин Артём Александрович}\\[0.2cm]
эксклюзивно для \textbf{Володи Хачатуряна}
}

\vspace{0.8cm}

{\large \today}

\vfill

\begin{tikzpicture}[x=1cm,y=1cm]
\draw[
  accent!65,
  line width=1.2pt
]
(0,0)
.. controls (2.3,0.55) and (3.8,-0.45) .. (6.0,0)
.. controls (8.1,0.43) and (10.1,-0.36) .. (12.0,0);
\end{tikzpicture}

\vspace{0.4cm}

\end{center}
\end{titlepage}

% ============================================================
% ОСНОВНОЙ ЧЕК-ЛИСТ
% ============================================================

\section{\color{darkaccent}1. Главная формула движения}
\sectionline

Для равномерного прямолинейного движения вдоль оси \(Ox\)
координата тела изменяется по закону

\[
\boxed{x=x_0+v_x t}.
\]

Здесь:

\begin{itemize}[leftmargin=*,itemsep=0.25em]
    \item \(x\) --- координата тела в момент времени \(t\);
    \item \(x_0\) --- начальная координата, то есть координата при \(t=0\);
    \item \(v_x\) --- проекция скорости на ось \(Ox\);
    \item \(t\) --- время, прошедшее от выбранного начального момента.
\end{itemize}

Если, например,

\[
x=-1-3t,
\]

то

\[
x_0=-1,\qquad v_x=-3.
\]

Через одну секунду:

\[
x(1)=-1-3\cdot1=-4.
\]

Так как \(v_x<0\), тело движется \important{против положительного направления оси \(Ox\)}.

% ------------------------------------------------------------

\section{\color{darkaccent}2. Связь с линейной функцией}
\sectionline

Линейная функция имеет вид

\[
\boxed{y=kx+b}.
\]

Её график --- прямая.

\begin{table}[h!]
\centering
\caption{Соответствие алгебры и физики}
\renewcommand{\arraystretch}{1.35}
\begin{tabularx}{\linewidth}{@{}XXX@{}}
\toprule
\textbf{Линейная функция} &
\textbf{График координаты} &
\textbf{Смысл} \\
\midrule
\(y\) & \(x\) & значение зависимой величины \\
\(x\) & \(t\) & независимая переменная \\
\(k\) & \(v_x\) & угловой коэффициент / скорость \\
\(b\) & \(x_0\) & свободный коэффициент / начальная координата \\
\bottomrule
\end{tabularx}
\end{table}

Следовательно,

\[
\boxed{x(t)=v_xt+x_0}
\]

--- это та же линейная зависимость, только обозначения имеют физический смысл.

\important{Важно:} буквы \(x\) и \(y\) сами по себе не определяют роль величины.
Сначала нужно установить, какая величина является независимой, а какая зависит от неё.

% ------------------------------------------------------------

\section{\color{darkaccent}3. Абсцисса и ордината}
\sectionline

На обычном графике функции:

\begin{itemize}[leftmargin=*,itemsep=0.25em]
    \item горизонтальная координата точки называется \important{абсциссой};
    \item вертикальная координата называется \important{ординатой}.
\end{itemize}

Точка

\[
A(2;5)
\]

имеет абсциссу \(2\) и ординату \(5\).

Для графика зависимости координаты от времени \(x(t)\):

\[
\boxed{\text{абсцисса}=t,\qquad \text{ордината}=x}.
\]

Поэтому точка

\[
A(2;5)
\]

на графике \(x(t)\) означает:

\[
t=2,\qquad x=5.
\]

\important{Одна из наиболее опасных ошибок} --- поменять \(t\) и \(x\) местами.

% ------------------------------------------------------------

\section{\color{darkaccent}4. Что показывает знак скорости}
\sectionline

В формуле

\[
x=x_0+v_xt
\]

коэффициент \(v_x\) выполняет роль углового коэффициента.

\begin{table}[h!]
\centering
\renewcommand{\arraystretch}{1.35}
\begin{tabularx}{\linewidth}{@{}XX@{}}
\toprule
\textbf{Условие} & \textbf{Физический смысл} \\
\midrule
\(v_x>0\) &
координата растёт; тело движется вдоль положительного направления \(Ox\) \\
\(v_x<0\) &
координата уменьшается; тело движется против положительного направления \(Ox\) \\
\(v_x=0\) &
координата постоянна; тело покоится \\
\bottomrule
\end{tabularx}
\end{table}

То же правило для \(y=kx+b\):

\[
k>0 \Rightarrow \text{функция возрастает},
\]

\[
k<0 \Rightarrow \text{функция убывает}.
\]

% ------------------------------------------------------------

\section{\color{darkaccent}5. Что показывает начальная координата}
\sectionline

При \(t=0\)

\[
x=x_0+v_x\cdot0=x_0.
\]

Поэтому

\[
\boxed{x_0=x(0)}.
\]

На графике \(x(t)\) начальная координата определяется по точке пересечения
графика с вертикальной осью \(x\).

Аналогично для линейной функции

\[
y=kx+b
\]

при \(x=0\)

\[
y=b.
\]

Следовательно, \(b\) --- ордината точки пересечения графика с осью \(Oy\).

% ------------------------------------------------------------

\section{\color{darkaccent}6. Четыре основные ситуации}
\sectionline

Знаки \(x_0\) и \(v_x\) отвечают за разные характеристики движения.

\begin{table}[h!]
\centering
\renewcommand{\arraystretch}{1.4}
\begin{tabularx}{\linewidth}{@{}ccX@{}}
\toprule
\(\boldsymbol{x_0}\) &
\(\boldsymbol{v_x}\) &
\textbf{Описание} \\
\midrule
\(>0\) & \(>0\) &
тело начинает справа от начала координат и движется дальше вдоль \(Ox\) \\
\(>0\) & \(<0\) &
тело начинает справа и движется в сторону начала координат \\
\(<0\) & \(>0\) &
тело начинает слева и движется в сторону начала координат \\
\(<0\) & \(<0\) &
тело начинает слева и движется дальше против \(Ox\) \\
\bottomrule
\end{tabularx}
\end{table}

Если

\[
x_0>0,\qquad v_x<0,
\]

то тело обязательно достигнет начала координат при некотором \(t>0\).

То же верно при

\[
x_0<0,\qquad v_x>0.
\]

% ------------------------------------------------------------

\section{\color{darkaccent}7. Когда тело проходит начало координат}
\sectionline

В начале координат

\[
x=0.
\]

Поэтому из

\[
x=x_0+v_xt
\]

получаем

\[
0=x_0+v_xt,
\]

откуда при \(v_x\neq0\)

\[
\boxed{t=-\frac{x_0}{v_x}}.
\]

Физически допустим момент

\[
\boxed{t\ge0}.
\]

Если вычисление даёт \(t<0\), это означает, что при рассматриваемом движении,
начиная с \(t=0\), тело через начало координат \important{не проходит}.

\subsection*{Пример}

Пусть

\[
x=5-2t.
\]

Тогда

\[
0=5-2t,
\]

\[
2t=5,
\]

\[
t=2{,}5.
\]

Через \(2{,}5\) с тело проходит начало координат.

% ------------------------------------------------------------

\section{\color{darkaccent}8. График координаты}
\sectionline

При построении графика \(x(t)\):

\begin{enumerate}[label=\arabic*),leftmargin=*,itemsep=0.3em]
    \item по горизонтали откладывайте время \(t\);
    \item по вертикали --- координату \(x\);
    \item обязательно подписывайте обе оси и единицы измерения;
    \item для обычной задачи движения рассматривайте \(t\ge0\);
    \item прямую проводите аккуратно по линейке;
    \item стрелки относятся к векторам, поэтому на самой линии графика их не ставят.
\end{enumerate}

\begin{center}
\begin{tikzpicture}
\begin{axis}[
    width=0.78\linewidth,
    height=7cm,
    axis lines=middle,
    unit vector ratio=1 1,
    xmin=0,
    xmax=7,
    ymin=-4,
    ymax=8,
    xlabel={$t$},
    ylabel={$x$},
    xtick={0,1,...,7},
    ytick={-4,-2,0,2,4,6,8},
    samples=220,
    clip=false,
    grid=both,
    minor tick num=0,
    grid style={line width=0.2pt,draw=gray!25},
    axis line style={draw=textgray},
    tick label style={font=\small},
    label style={font=\small}
]
\addplot[
    domain=0:6.2,
    samples=220,
    line width=1.4pt,
    color=accent
]{3+0.7*x};
\addplot[
    domain=0:6.2,
    samples=220,
    line width=1.4pt,
    color=darkaccent,
    dashed
]{5-x};

\addplot[
    only marks,
    mark=*,
    mark size=2pt,
    color=accent
] coordinates {(0,3)};

\addplot[
    only marks,
    mark=*,
    mark size=2pt,
    color=darkaccent
] coordinates {(0,5) (5,0)};

\node[anchor=west,font=\small] at (axis cs:4.3,6.2)
{$v_x>0$};

\node[anchor=west,font=\small] at (axis cs:3.1,1.6)
{$v_x<0$};
\end{axis}
\end{tikzpicture}
\end{center}

У сплошного графика координата возрастает, следовательно \(v_x>0\).

У пунктирного графика координата уменьшается, поэтому \(v_x<0\).
Он пересекает ось времени в точке \(t=5\), следовательно в этот момент \(x=0\).

% ------------------------------------------------------------

\section{\color{darkaccent}9. Как найти скорость по графику}
\sectionline

Для равномерного движения скорость равна отношению изменения координаты
к соответствующему промежутку времени:

\[
\boxed{
v_x=\frac{\Delta x}{\Delta t}
=
\frac{x_2-x_1}{t_2-t_1}
}.
\]

Обязательно:

\[
t_2\neq t_1.
\]

Если координата измеряется в метрах, а время в секундах, то скорость
получается в \(\text{м/с}\).

\subsection*{Пример}

На графике отмечены две точки:

\[
A(2;1),\qquad B(6;2),
\]

где первая координата --- \(t\), вторая --- \(x\).

Тогда

\[
v_x=\frac{2-1}{6-2}
=\frac14.
\]

Начальную координату найдём из

\[
x=x_0+v_xt.
\]

Подставим точку \(A\):

\[
1=x_0+\frac14\cdot2,
\]

\[
1=x_0+\frac12,
\]

\[
\boxed{x_0=\frac12}.
\]

Следовательно,

\[
\boxed{x=\frac12+\frac14t}.
\]

% ------------------------------------------------------------

\section{\color{darkaccent}10. Метод двух узловых точек}
\sectionline

В задачах этого типа удобно выбирать две точки графика, координаты которых
легко и точно считываются по сетке. Такие точки часто называют
\important{узловыми точками}.

Для линейной функции

\[
y=kx+b
\]

неизвестны два коэффициента: \(k\) и \(b\).

Поэтому достаточно двух различных точек графика.

Пусть

\[
A(0;-1),\qquad B(2;0).
\]

Подставляем координаты в формулу:

\[
\begin{cases}
-1=k\cdot0+b,\\
0=k\cdot2+b.
\end{cases}
\]

Получаем

\[
\begin{cases}
b=-1,\\
2k+b=0.
\end{cases}
\]

Подставляем \(b=-1\):

\[
2k-1=0,
\]

\[
2k=1,
\]

\[
k=\frac12.
\]

Итак,

\[
\boxed{y=\frac12x-1}.
\]

\begin{center}
\begin{tikzpicture}
\begin{axis}[
    width=0.72\linewidth,
    height=6.6cm,
    axis lines=middle,
    unit vector ratio=1 1,
    xmin=-1,
    xmax=6,
    ymin=-3,
    ymax=3,
    xlabel={$x$},
    ylabel={$y$},
    xtick={-1,0,...,6},
    ytick={-3,-2,...,3},
    samples=220,
    clip=false,
    grid=both,
    grid style={line width=0.2pt,draw=gray!25},
    axis line style={draw=textgray},
    tick label style={font=\small},
    label style={font=\small}
]
\addplot[
    domain=-1:6,
    samples=220,
    line width=1.4pt,
    color=accent
]{0.5*x-1};

\addplot[
    only marks,
    mark=*,
    mark size=2.4pt,
    color=darkaccent
] coordinates {(0,-1) (2,0) (4,1)};

\node[anchor=east,font=\small] at (axis cs:0,-1) {$A$};
\node[anchor=north west,font=\small] at (axis cs:2,0) {$B$};
\node[anchor=south east,font=\small] at (axis cs:4,1) {$C$};
\end{axis}
\end{tikzpicture}
\end{center}

% ------------------------------------------------------------

\section{\color{darkaccent}11. Быстрый способ через две точки}
\sectionline

Если известны две точки прямой

\[
(x_1;y_1),\qquad (x_2;y_2),
\]

то

\[
\boxed{k=\frac{y_2-y_1}{x_2-x_1}},
\qquad x_2\neq x_1.
\]

После этого

\[
\boxed{b=y_1-kx_1}.
\]

Для графика координаты используются полностью аналогичные формулы:

\[
\boxed{
v_x=\frac{x_2-x_1}{t_2-t_1}
},
\]

\[
\boxed{x_0=x_1-v_xt_1}.
\]

Это тот же метод, записанный в более компактной форме.

% ------------------------------------------------------------

\section{\color{darkaccent}12. Смысловая проверка ответа}
\sectionline

После вычислений обязательно сопоставьте ответ с графиком.

Например, если график координаты явно возрастает, а получилось

\[
v_x=0
\]

или

\[
v_x<0,
\]

то вычисления содержат ошибку.

Полезная проверка:

\begin{itemize}[leftmargin=*,itemsep=0.35em]
    \item график возрастает \(\Rightarrow v_x>0\);
    \item график убывает \(\Rightarrow v_x<0\);
    \item горизонтальный график \(\Rightarrow v_x=0\);
    \item пересечение вертикальной оси должно давать \(x_0\);
    \item при подстановке любой точки графика равенство
    \(x=x_0+v_xt\) должно выполняться.
\end{itemize}

% ------------------------------------------------------------

\section{\color{darkaccent}13. Алгебра: как правильно избавляться от коэффициентов}
\sectionline

При решении уравнения выполняйте \important{одно и то же действие с обеими частями}.

Например,

\[
4v_x=1.
\]

Чтобы оставить слева только \(v_x\), нужно разделить обе части на \(4\):

\[
\frac{4v_x}{4}=\frac14,
\]

поэтому

\[
\boxed{v_x=\frac14}.
\]

Знак минус здесь не возникает.

Другой пример:

\[
5v_x+5=0.
\]

Сначала вычтем \(5\) из обеих частей:

\[
5v_x=-5.
\]

Теперь разделим на \(5\):

\[
\boxed{v_x=-1}.
\]

Фраза «перенести через знак равенства» является лишь сокращённым способом
описать эквивалентные преобразования уравнения.

% ------------------------------------------------------------

\section{\color{darkaccent}14. Типичные ошибки}
\sectionline

\begin{enumerate}[label=\arabic*),leftmargin=*,itemsep=0.55em]

    \item \textbf{Перепутаны координаты точки.}

    Для графика \(x(t)\)

    \[
    A(2;5)
    \]

    означает \(t=2\), \(x=5\).

    \item \textbf{Знак скорости определён по положению графика, а не по его наклону.}

    График может находиться выше оси времени и при этом убывать.

    \item \textbf{Начальная координата перепутана со скоростью.}

    В

    \[
    x=x_0+v_xt
    \]

    на \(t\) умножается именно \(v_x\).

    \item \textbf{При решении \(4v_x=1\) появляется лишний минус.}

    Правильно:

    \[
    v_x=\frac14.
    \]

    \item \textbf{Не проведена смысловая проверка.}

    Если вычисленный знак скорости противоречит наклону графика,
    нужно проверить арифметику и порядок координат.

    \item \textbf{Не подписаны оси.}

    Всегда указывайте, где \(t\), где \(x\), а при необходимости ---
    секунды, метры и другие единицы.

    \item \textbf{Для времени принят отрицательный ответ.}

    Если движение рассматривается начиная с \(t=0\), физическая область:

    \[
    t\ge0.
    \]

    \item \textbf{График построен неаккуратно.}

    Прямые линии проводите по линейке.

\end{enumerate}

% ------------------------------------------------------------

\section{\color{darkaccent}15. Итоговый алгоритм}
\sectionline

Если дан график зависимости координаты от времени и требуется описать движение:

\begin{enumerate}[label=\arabic*),leftmargin=*,itemsep=0.45em]

    \item Проверьте подписи осей:
    \[
    t \text{ --- горизонталь},\qquad x \text{ --- вертикаль}.
    \]

    \item Найдите начальную координату:
    \[
    x_0=x(0).
    \]

    \item Определите знак скорости по направлению графика:
    \[
    \text{возрастает} \Rightarrow v_x>0,
    \]
    \[
    \text{убывает} \Rightarrow v_x<0.
    \]

    \item Выберите две точно считываемые точки.

    \item Рассчитайте скорость:
    \[
    v_x=\frac{x_2-x_1}{t_2-t_1}.
    \]

    \item Рассчитайте начальную координату:
    \[
    x_0=x_1-v_xt_1.
    \]

    \item Запишите закон движения:
    \[
    x=x_0+v_xt.
    \]

    \item Подставьте одну из точек обратно и проверьте равенство.

    \item Сопоставьте полученный знак скорости с наклоном графика.

    \item Если нужно найти прохождение начала координат, решите:
    \[
    0=x_0+v_xt
    \]
    и проверьте условие \(t\ge0\).

\end{enumerate}

% ------------------------------------------------------------

\section{\color{darkaccent}16. Минимум, который нужно знать без подсказки}
\sectionline

\begin{itemize}[leftmargin=*,itemsep=0.5em]
    \item[\checkitem] Формулу
    \[
    x=x_0+v_xt.
    \]

    \item[\checkitem] Связь
    \[
    x=x_0+v_xt
    \quad\longleftrightarrow\quad
    y=kx+b.
    \]

    \item[\checkitem] На графике \(x(t)\):
    \[
    \text{горизонталь}=t,\qquad
    \text{вертикаль}=x.
    \]

    \item[\checkitem]
    \[
    x_0=x(0).
    \]

    \item[\checkitem]
    \[
    v_x>0 \Rightarrow x \text{ растёт},
    \qquad
    v_x<0 \Rightarrow x \text{ уменьшается}.
    \]

    \item[\checkitem]
    \[
    v_x=\frac{x_2-x_1}{t_2-t_1}.
    \]

    \item[\checkitem]
    \[
    t_{\text{при }x=0}=-\frac{x_0}{v_x},
    \qquad v_x\neq0.
    \]

    \item[\checkitem]
    Физически допустимое время после начала наблюдения:
    \[
    t\ge0.
    \]

    \item[\checkitem]
    После вычислений ответ обязательно проверяется по смыслу графика.
\end{itemize}

% ============================================================
% ТРЕНИРОВОЧНЫЙ БЛОК
% ============================================================

\clearpage

\section{\color{darkaccent}Тренировочный блок}
\sectionline

\textbf{Путь Б.} Выполняйте по пять заданий в день.
Во всех физических задачах время \(t\) измеряется в секундах,
координата \(x\) --- в метрах, скорость --- в \(\text{м/с}\),
если отдельно не указано иное.

% ------------------------------------------------------------

\daytitle{06.09}

\begin{enumerate}[label=\arabic*),leftmargin=*,itemsep=0.9em]

\item Движение задано уравнением
\[
x=-3+2t.
\]
Найдите \(x_0\), \(v_x\), координату через \(4\) с и укажите направление движения.

\item Движение задано уравнением
\[
x=5-1{,}5t.
\]
Через какое время тело пройдёт начало координат?

\item На графике \(x(t)\) выбраны точки
\[
A(2;1),\qquad B(6;2).
\]
Найдите \(v_x\), \(x_0\) и запишите уравнение движения.

\item Для функции
\[
y=-2x+3
\]
найдите \(k\) и \(b\). Укажите, возрастает или убывает функция и в какой точке график пересекает ось \(Oy\).

\item Известно:
\[
x_0>0,\qquad v_x<0.
\]
Опишите расположение и направление графика \(x(t)\). Обязательно ли тело при \(t>0\) пройдёт начало координат?

\end{enumerate}

% ------------------------------------------------------------

\daytitle{07.09}

\begin{enumerate}[label=\arabic*),leftmargin=*,itemsep=0.9em]

\item Движение задано уравнением
\[
x=-2-3t.
\]
Найдите \(x_0\), \(v_x\), \(x(2)\) и направление движения.

\item Дано:
\[
x=6+1{,}5t.
\]
В какой момент координата станет равной \(12\)?

\item На графике \(x(t)\) выбраны точки
\[
A(1;4),\qquad B(5;0).
\]
Найдите скорость, начальную координату и закон движения.

\item Для функции
\[
y=\frac34x-2
\]
определите \(k\), \(b\), характер монотонности и ординату точки пересечения с \(Oy\).

\item Тело имеет
\[
x_0=-4,\qquad v_x=0{,}5.
\]
Запишите закон движения и найдите момент прохождения начала координат.

\end{enumerate}

% ------------------------------------------------------------

\daytitle{08.09}

\begin{enumerate}[label=\arabic*),leftmargin=*,itemsep=0.9em]

\item Дано:
\[
x=7-2t.
\]
Найдите \(x(3)\) и момент прохождения начала координат.

\item На графике \(x(t)\) имеются точки
\[
A(0;-4),\qquad B(3;2).
\]
Найдите \(x_0\), \(v_x\) и уравнение движения.

\item Рассмотрите функцию
\[
y=-\frac12x-3.
\]
Найдите \(k\), \(b\), укажите характер монотонности и знак свободного коэффициента.

\item Движение описывается формулой
\[
x=-5+2{,}5t.
\]
Через сколько секунд координата станет равной \(10\)?

\item На графике \(x(t)\) выбраны точки
\[
A(2;-1),\qquad B(5;8).
\]
Найдите \(v_x\), \(x_0\) и уравнение движения.

\end{enumerate}

% ------------------------------------------------------------

\daytitle{09.09}

\begin{enumerate}[label=\arabic*),leftmargin=*,itemsep=0.9em]

\item Движение задано уравнением
\[
x=1+\frac12t.
\]
Найдите начальную координату, скорость и координату через \(6\) с.

\item На графике \(x(t)\) выбраны точки
\[
A(4;3),\qquad B(10;-3).
\]
Найдите \(v_x\), \(x_0\) и закон движения.

\item Для функции
\[
y=-4x+6
\]
найдите \(k\), \(b\). Определите, возрастает или убывает функция.

\item Известно:
\[
x_0=-6,\qquad v_x=2.
\]
Запишите уравнение движения и найдите момент прохождения начала координат.

\item Тело имеет
\[
x_0=9,\qquad v_x=-3.
\]
Найдите координату через \(5\) с и момент прохождения начала координат.

\end{enumerate}

% ------------------------------------------------------------

\daytitle{10.09}

\begin{enumerate}[label=\arabic*),leftmargin=*,itemsep=0.9em]

\item Дано:
\[
x=-1+\frac32t.
\]
Найдите \(x(4)\) и момент прохождения начала координат.

\item На графике \(x(t)\) выбраны точки
\[
A(2;5),\qquad B(8;2).
\]
Найдите скорость, начальную координату и закон движения.

\item Для функции
\[
y=2{,}5x+1
\]
найдите \(k\), \(b\), определите характер монотонности и точку пересечения с осью \(Oy\).

\item Рассмотрите движение
\[
x=4-0{,}8t.
\]
Укажите знаки \(x_0\) и \(v_x\), направление движения и момент прохождения начала координат.

\item Ученик решил уравнение
\[
4v_x=1
\]
и записал
\[
v_x=-4.
\]
Исправьте ошибку и запишите правильное значение \(v_x\).

\end{enumerate}

% ------------------------------------------------------------

\daytitle{11.09}

\begin{enumerate}[label=\arabic*),leftmargin=*,itemsep=0.9em]

\item Движение задано уравнением
\[
x=8-2{,}5t.
\]
Найдите \(x(2)\) и момент прохождения начала координат.

\item На графике \(x(t)\) выбраны точки
\[
A(3;-2),\qquad B(7;6).
\]
Найдите \(v_x\), \(x_0\) и уравнение движения.

\item Для функции
\[
y=-\frac13x+4
\]
определите \(k\), \(b\), характер монотонности и ординату точки пересечения с \(Oy\).

\item Известно:
\[
x_0=-3,\qquad v_x=-2.
\]
Запишите закон движения, найдите \(x(4)\) и определите, проходит ли тело начало координат при \(t\ge0\).

\item На графике координаты отмечены точки
\[
A(1;7),\qquad B(5;-1).
\]
Найдите скорость, начальную координату, закон движения и момент прохождения начала координат.

\end{enumerate}

% ============================================================
% ОТВЕТЫ
% ============================================================

\clearpage

\section{\color{darkaccent}Ответы к тренировочному блоку}
\sectionline

\renewcommand{\arraystretch}{1.35}
\small

\begin{longtable}{@{}p{1.6cm}p{0.8cm}p{\dimexpr\linewidth-3.0cm\relax}@{}}
\toprule
\textbf{Дата} & \textbf{№} & \textbf{Ответ} \\
\midrule
\endfirsthead

\toprule
\textbf{Дата} & \textbf{№} & \textbf{Ответ} \\
\midrule
\endhead

\midrule
\multicolumn{3}{r}{\small Продолжение на следующей странице} \\
\endfoot

\bottomrule
\endlastfoot

06.09 & 1 &
\(x_0=-3\), \(v_x=2\), \(x(4)=5\); движение вдоль положительного направления \(Ox\). \\

06.09 & 2 &
\(\displaystyle t=\frac{10}{3}\text{ с}\approx3{,}33\text{ с}\). \\

06.09 & 3 &
\(\displaystyle v_x=\frac14\), \(\displaystyle x_0=\frac12\);
\(\displaystyle x=\frac12+\frac14t\). \\

06.09 & 4 &
\(k=-2\), \(b=3\); функция убывает; пересечение с \(Oy\): \((0;3)\). \\

06.09 & 5 &
График начинается выше оси \(t\) и убывает. Да:
\(\displaystyle t=-\frac{x_0}{v_x}>0\). \\

\midrule

07.09 & 1 &
\(x_0=-2\), \(v_x=-3\), \(x(2)=-8\); движение против положительного направления \(Ox\). \\

07.09 & 2 &
\(t=4\text{ с}\). \\

07.09 & 3 &
\(v_x=-1\), \(x_0=5\);
\[
x=5-t.
\] \\

07.09 & 4 &
\(\displaystyle k=\frac34\), \(b=-2\); функция возрастает; пересечение с \(Oy\): \((0;-2)\). \\

07.09 & 5 &
\[
x=-4+0{,}5t,\qquad t=8\text{ с}.
\] \\

\midrule

08.09 & 1 &
\(x(3)=1\);
\(\displaystyle t=\frac72=3{,}5\text{ с}\). \\

08.09 & 2 &
\(x_0=-4\), \(v_x=2\);
\[
x=-4+2t.
\] \\

08.09 & 3 &
\(\displaystyle k=-\frac12\), \(b=-3\); функция убывает; \(b<0\). \\

08.09 & 4 &
\(t=6\text{ с}\). \\

08.09 & 5 &
\(v_x=3\), \(x_0=-7\);
\[
x=-7+3t.
\] \\

\midrule

09.09 & 1 &
\(\displaystyle x_0=1\), \(\displaystyle v_x=\frac12\), \(x(6)=4\). \\

09.09 & 2 &
\(v_x=-1\), \(x_0=7\);
\[
x=7-t.
\] \\

09.09 & 3 &
\(k=-4\), \(b=6\); функция убывает. \\

09.09 & 4 &
\[
x=-6+2t,\qquad t=3\text{ с}.
\] \\

09.09 & 5 &
\(x(5)=-6\);
\(t=3\text{ с}\). \\

\midrule

10.09 & 1 &
\(x(4)=5\);
\(\displaystyle t=\frac23\text{ с}\). \\

10.09 & 2 &
\(\displaystyle v_x=-\frac12\), \(x_0=6\);
\[
x=6-\frac12t.
\] \\

10.09 & 3 &
\(k=2{,}5\), \(b=1\); функция возрастает; пересечение с \(Oy\): \((0;1)\). \\

10.09 & 4 &
\(x_0>0\), \(v_x<0\); движение против положительного направления \(Ox\);
\(t=5\text{ с}\). \\

10.09 & 5 &
\[
\boxed{v_x=\frac14}.
\]
Нужно разделить обе части уравнения на \(4\). \\

\midrule

11.09 & 1 &
\(x(2)=3\);
\(\displaystyle t=\frac{16}{5}=3{,}2\text{ с}\). \\

11.09 & 2 &
\(v_x=2\), \(x_0=-8\);
\[
x=-8+2t.
\] \\

11.09 & 3 &
\(\displaystyle k=-\frac13\), \(b=4\); функция убывает; пересечение с \(Oy\): \((0;4)\). \\

11.09 & 4 &
\[
x=-3-2t,\qquad x(4)=-11.
\]
Из \(0=-3-2t\) получается \(t=-1{,}5\), поэтому при \(t\ge0\) тело начало координат не проходит. \\

11.09 & 5 &
\(v_x=-2\), \(x_0=9\);
\[
x=9-2t,
\qquad
t=\frac92=4{,}5\text{ с}.
\] \\

\end{longtable}

\normalsize

\end{document}